\begin{abstract}
La reducción dimensional modifica las vecindades de las representaciones contextuales. Estudiamos si los vecinos conservados o introducidos por PCA, ISOMAP y UMAP comparten el synset de WordNet en \Num{n} ocurrencias de SemCor distribuidas en \Num{groups} grupos lema–POS. Reducimos al plano embeddings BERT normalizados y comparamos retención ajustada por azar, pureza y selectividad semántica. Una identidad exacta separa el cambio de pureza en fracción de vecinos sustituidos y diferencia de pureza entre los vecinos introducidos y eliminados. En la configuración principal, UMAP obtiene la mayor AUC macro de conservación en una ventana con pesos inversos al tamaño de la vecindad. Las bandas simultáneas de sus dos contrastes de conservación excluyen cero. La diferencia es mayor para vecindades pequeñas y se invierte a $k=9$ en la configuración principal. UMAP también presenta el mayor cambio puntual de pureza ajustada en la configuración principal. Las bandas de los contrastes semánticos incluyen cero. El cambio de pureza ajustada de UMAP varía de signo según su escala. La retención de identidades y la coincidencia de synsets son criterios lógicamente distintos. Las bandas son aproximaciones condicionales a la geometría fija y a la reponderación de los documentos observados. Los resultados se limitan a un piloto exploratorio.
\end{abstract}
\noindent\textbf{Palabras clave:} polisemia, embeddings contextuales, reducción dimensional, conservación de vecindarios, WordNet.
